\documentclass[eprint,
%approc,		%heading for Acta Polytechnica CTU Proceedings
]{actapoly}

\usepackage[utf8]{inputenc}
\usepackage{soul}
\usepackage{caption}
\usepackage{subcaption}
\usepackage{array}
\usepackage[flushleft]{threeparttable}

\newcommand*\printft[1]{\textsc{\lowercase{#1}}}

\graphicspath{{figures/}} % path to folder with figures
\usepackage[super]{nth}
\usepackage{fixltx2e}
\usepackage{tabularx,booktabs}
	\newcolumntype{Y}{>{\centering\arraybackslash}X}
\usepackage{amsmath}
\usepackage{mathtools, nccmath}

\begin{document}

\title[Pareto Efficiency and RSM to Adjust Binder Content]
{Utilising Pareto Efficiency and RSM to Adjust Binder Content in Clay Stabilisation for Yttre Ringvägen, Malmö}

\author[P. Lindh]{Per Lindh}{my,their}
\correspondingauthor[P. Lemenkova]{Polina Lemenkova}{third}{polina.lemenkova@ulb.be}

\institution{my}{Swedish Transport Administration, Department of Investments Technology and Environment, Neptunigatan 52, Box 366, SE-201-23 Malmö, Sweden}
\institution{their}{Lund University, Lunds Tekniska Högskola LTH (Faculty of Engineering), Department of Building and Environmental Technology, Division of Building Materials, Box 118, SE-221-00, Lund, Sweden}
\institution{third}{Université Libre de Bruxelles (ULB), École polytechnique de Bruxelles (Brussels Faculty of Engineering), Laboratory of Image Synthesis and Analysis (LISA). Campus de Solbosch, ULB -- LISA CP165/57, Avenue Franklin D. Roosevelt 50, B-1050 Brussels, Belgium}

\begin{abstract}
In this paper, we present a new framework for improving soil strength using an advanced method of engineering statistics. The materials included clay till collected in Yttre Ringvägen, southern Sweden. Binders included quicklime, slag and ordinary Portland cement used as pure binders and blended mixtures. We first applied the Response Surface Methodology techniques aimed at binder blend optimisation: 1) Central Composite Design; 2) Box-Behnken Design; 3) Simplex Lattice Design. The Pareto charts were presented for modelling responses from tests with different binders and estimating their effects on soil strength. Finally, to examine the variables important for soil stabilisation, we also evaluated the effect of the amount of binder and the interaction between cement/lime/slag in different ratios: 30-50-20\%; 50-50-0\%; 100-0-0\% The paper highlights the major opportunities and challenges of engineering statistics as a cross-cutting research direction for the issues of civil engineering.
\end{abstract}

\keywords{soil stabilisation, simplex experimental design, binder, OPC, statistical analysis}

\maketitle

\section{Introduction}

Soil \hl{stabilisation} is a critically important task in civil engineering. It is aimed at improving soil parameters and properties in various areas of civil engineering\hl{, }such as road constructions, bridge or building engineering, \hl{and} earthworks on pavements. \hl{Stabilisation} and solidification of soil is a widely applied method in geotechnical works performed using various binders as stabilising agents. The existing methods of soil \hl{stabilisation are} aimed at improving \hl{the} soil performance to obtain the required characteristics of foundations. Although the state-of-the art methods are applicable, new approaches using various stabilisation agents used \hl{solitarily} or in combinations and novel binders require experimental testing and evaluation to assess their quality and effectiveness. 

Despite widespread applications in civil engineering and extensive existing research, there are still some challenges in soil \hl{stabilisation,} including the following: 1) Evaluation of variations in soil properties and responses of individual soil specimens and variations between binder characteristics that affect \hl{the stabilisation} process. For instance, this includes \hl{the} ratio of \hl{stabilising} agents with respect to water content, technical characteristics of \hl{the stabilising} agent and the use of accelerators; 2) Processing large \hl{amounts} of sampling materials – modern engineering tasks in construction industry are increasingly high dimensional and require\hl{, }in many cases\hl{, }processing of several tons of soil using hundreds of kg of binders. This necessarily requires the use of the effective techniques to optimise this process; 3) Dealing with occlusions in soil such as fiber, which may create noise in technical parameters of soil while modelling data\hl{, }and therefore affect \hl{the stabilisation} process. 

At the same time, effective soil \hl{stabilisation} is crucial for safe road constructions and engineering works. This especially concerns northern regions with harsh environmental and climate conditions that create challenges for infrastructure \citep{Dahlinetal1999,ROTHHAMEL2022100735,Lindh2004,VESTIN2018146,Lemenkov2021a,Lemenkov2021b,Eigenbrod}\hl{, due to the} unique physical and mechanical properties of soil collected in real-world environment, rather than theoretical models presented in technical guidance. \hl{Taking this into account}, the importance of experiments on soil \hl{stabilisation} consists in the complexity of real case situations because soil is a highly variable porous structure. Formed as a mixture of organic matter, minerals and rocks, chemical components (gases or liquids), and biological particles (microorganisms)\hl{, }the properties of soil vary significantly\hl{, }which requires experimental testing in \hl{earthwork} constructions. 

While much progress has been made in soil \hl{stabilisation} techniques to address various aspects of the \hl{abovementioned} challenges, many issues are still far from being resolved. Particularly, \hl{those that} lack an effective and standardised framework for considering binder optimisation during the process of soil stabilization. Thus, the majority of the accepted \emph{in situ} techniques of soil \hl{stabilisation} are based on the approach \hl{where the stabilisation} agent\hl{, }as one binder\hl{, }is distributed over soil specimens. In such cases, the most widely used binder is lime, as proved by numerous examples of the existing cases \citep{Greaves,RaoThyagaraj,Rogersetal2000,ATTOHOKINE1995283,ROBIN2014211,Harichane,GlendinningRogers}. The Ordinary Portland Cement (OPC) is another widely used \hl{stabilising} agent. Recent advances in utilizing OPC as a binder \citep{Prusinski,Bhattacharjaetal2003,Lindh202201,Lindh202102,KULKARNI2022125363} have made the OPC \hl{widely used} for \hl{stabilisation} of soil in the field of civil engineering. The difference in the effects of OPC and other binders poses\hl{ }a new way to model and compare the reaction of soil with binders during the \hl{stabilisation} process. For instance, the reactions of lime and OPC with soil differ and have \hl{their own} advantages and disadvantages, as discussed in existing literature \citep{Mahedi,KOLIAS2005301,BARMAN20221319,Lindhetal2000,Chewetal2004,KOOHMISHI2022100726}. 

The advantages of lime\hl{, }which is one of the oldest binders used to improve engineering properties of soils, are as follows: 1) it results in\hl{ }long period of exploitation and workability; 2) it enables high level of homogeneity in the mixtures with soil; 3) it is effective \hl{in decreasing} the amount of water; 4) it reduces the plasticity index of soil, which facilitates higher workability of foundations \cite{LindhLemenkovaCEE}. However, there are also certain disadvantages of lime\hl{, }which can be mentioned as follows: 1) lime needs certain conditions regarding the mineral content and grading of soil; 2) although lime does increase the strength of soil, the process is rather slow; 3) organic content of soil may affect the performance of lime as a binder.

The advantages of OPC as a binder include the following ones: 1) OPC ensures \hl{a} high level of strength during the process of \hl{stabilisation}; 2) using OPC increases the speed of the strength gain; 3) OPC \hl{reacts well} with water and thus enables fabricating the OPC slurry; 4) Compared to lime, OPC has higher robustness regarding soil grading; 5) organic content in soil does not affect the OPC\hl{'s} performance as much \hl{in the case of lime}. Nevertheless, some drawbacks of the OPC as a binder should also be mentioned: 1) the effects \hl{resulting from the properties of} OPC remain for a shorter time; 2) the use of OPC \hl{requires a} significant amount of compaction works; 3) using OPC results in \hl{a} lower homogeneity of \hl{the} OPC-soil mixture\hl{ as} compared to lime-soil mixture.

The performance of \hl{the} binder during soil \hl{stabilisation} and its reaction with specimens ultimately explains the difference in the effects \hl{of} lime and OPC, caused by \hl{the} mineral structure. Thus, aluminium and silica minerals and water are necessary for lime to produce pozzolanic reactions. Their main effects include soil cementation with\hl{ }higher strength, reduced deformability\hl{,} and higher durability \citep{DISANTE2022100742,AkulaLittle,1988240,VITALE201736}. These minerals can already be present naturally or can be mixed with lime during \hl{the} soil \hl{stabilisation}. The advantage of the pozzolanic reaction is that it ensures the increase in strength of \hl{soil} which lasts over years. In contrast, the OPC or cementitious binders are less sensitive and only need water for \hl{an} effective reaction. In this case, hardening starts immediately and the gain of strength in the OPC-soil mixture increases exponentially with the maximum value achieved after 28 days of curing time \citep{Bell1995,Makusa}. 

Besides the effects from binders, various soil types behave differently during the \hl{stabilisation} process \citep{BhattacharjaBhatty,Burroughs2008,Renjithetal2020,VENDAOLIVEIRA2018336,Wall}. Therefore, there is no unique recipe for the mixture of stabilisation agents. As a response to this problem, the objective of the \hl{stabilisation} methods is to correctly select binders and adjust their amounts and ratios in order to achieve the most effective \hl{stabilisation} results. Hence, binders should be defined and regulated carefully with respect \hl{to} the \emph{in situ} conditions and parameters of soil. For example, blended binders may sometimes ensure the best performance, while in other cases, the effects from various binders on soil should be tested empirically. Very often, the use of blended binders results in the best performance of the binder-soil mixture\hl{,} which is the goal of \hl{stabilisation} and required workability of \hl{soil} \citep{Schanz,Kangetal2019}. 

In this paper, we propose a framework for a number of\hl{ }experimental designs for establishing an optimal technique of soli \hl{stabilisation} tested for different types and blended binders. The objective of the study is to contribute to the development of technical methods of robust soil \hl{stabilisation} through empirical binder optimisation using a combination of the statistical approach and series of practical tests. The goal was to find and indicate the optimal mixture of \hl{stabilising} agents for soil \hl{stabilisation} considering both \hl{the} soil properties and \hl{the} economical limitations. The experiment revealed both the effective and not effective interactions between the different \hl{stabilising} agents and their reaction with soil specimens over time during the \hl{stabilisation} process. 

The project is based on the empirical geotechnical works performed in the laboratory of the Swedish Geotechnical institute (SGI). The technical aim was to fabricate blended mixtures made from different pure binders \hl{having the best effect on soil stabilisation}. To this end, we used OPC, lime, GGBFS and fly ash as binders for \hl{stabilising} soil samples. The framework included a number of existing standardised technical \hl{workflows} of soil \hl{stabilisation, }modified and applied for our case.

\section{Essentials on binder selection}
A mixture of two or three different binders is commonly acceptable for soil \hl{stabilisation} in general and for Deep Mixing Method (DMM) in particular \cite{BergmanLarsson,Naqshabandy,Paniagua}. In the last case, the lime-OPC columns are fabricated using a special machine. Previous research shown that a mixture of OPC and fly ash contributes to \hl{the} gain of soil strength better \hl{as} compared to OPC \hl{used} as a single binder \cite{Yaxuetal2022,Vukicevicetal}. Other cases shown that there is a positive reaction between lime and fly ash\hl{, }which can be used to improve the results of \hl{stabilisation} \cite{GoswamiSingh,Raoetal2008,Knapiketal2015}. 

Blended binders have been successfully used for the past 30 years for the purpose of soil stabilization \cite{Ghazy,Consolietal2019,Lindh2001,Leeetal2019,Filho,Yietal2014,LindhLemenkovaECES,Guetal2014,Yangetal2015}. To improve the efficiency of soil stabilisation, some commercial binder suppliers developed their own blends. Such binders can be classified as follows: 1). General Purpose (GP) OPC; 2). General Blend (GB) OPC; 3). Cementitious triple and quaternary blends of binders mixed from combinations of fly ash, GP OPC, ground granulated blast furnace slag (GGBFS) and lime; 4). Hydrated lime; 5). OPC/asphalt blends. Previous studies on the evaluation of the effects \hl{of} various \hl{mixtures} demonstrated \cite{ma15217798} that GGBFS contributes the most to the compressive strength development in expansive clayey soil, followed by OPC. Moreover, blended mixtures of cement with GGBFS and lime with GGBFS perform better than single binders. Therefore, we chose a GGBS combined with cement and lime as \hl{the} mixture to achieve notable effects on soil \hl{stabilisation} and strength gain. In Sweden, the use of blended binders is \hl{on the rise} since 1970s, Figure \ref{fig1}. The Swedish OPC and lime industry encouraged the use of the additives with regard to soil type, Figure \ref{fig2}. Nowadays, the use of mixed binders in deep mixing have increased up to 100\%, see Figure \ref{fig1}. The applicability of different binders is summarised in Table \ref{tab:1}. 

%\begin{table}[htbp]
\begin{table*}[t]
\begin{threeparttable}
\centering
\caption{Applicability of binders for soil types. Source: modified after \cite{AustStab} \label{tab:1}}
\begin{tabularx}{\textwidth}{c *{6}{Y}}
\toprule
Binder type & \multicolumn{4}{c}{Course-grained soil} & \multicolumn{2}{c}{Fine-grained soil}\\
\cmidrule(lr){2-5} \cmidrule(l){6-7}
  & GP & GW & GM \& GC & SP \& SW & CL & CH \\
\midrule
GP OPC (A) & 1 & 1 & 1 & 2 & 2 & 3 \\ 
GB OPC (B) & 1 & 1 & 1 & 1 & 1 & 2 \\ 
OPCitious blends (C) &  1 & 1 & 1 & 1 & 1 & 3 \\ 
Lime \& OPC (C) & 3 & 3 & 2 & 3 & 2 & 1 \\ 
Lime \& fly ash (C) & 3 & 1 & 1 & 3 & 2 & 2 \\ 
Lime (D) & 2 & 2 & 1 & 3 & 2 & 1 \\ 
OPC/bitumen (E) & 1 & 1 & 2 & 2 & 3 & 3 \\ 
\bottomrule
\end{tabularx}
\begin{tablenotes}
      \small
      	\item \emph{Notations for Table 1.} GP - Poorly graded gravel and \hl{crushed} rock; GW - Well-graded gravel; GM - Silty gravel; GC - Clayey gravel; SP - clean sand, poorly graded; SW - Clean sand - well graded; CL - Low plasticity sandy/silty clay; CH - High plasticity heavy clay. 1 – Excellent; 2 – Satisfactory; 3 – Not suitable.
    \end{tablenotes}
  \end{threeparttable}
\end{table*}

However, there is \hl{a} lack of \hl{documented} experience regarding the optimisation of different mixtures of binders for stabilisation purposes\hl{, }specifically for Sweden. Selected papers reported cases of binder optimisation using deep mixing methods, for instance\hl{, }at the SGI \cite{Holm2001,Holm2003}. Regarding the strength development under the effects of additives added to different soil \hl{types}, \hl{the SGI documented the following: the} OPC usually demonstrates very good effects for high-plasticity clayey silt (MH), lime-OPC blend gives \hl{a} good effect, \hl{and} lime is satisfactory for soil \hl{stabilisation}. 

\begin{figure}[H]\centering
	\includegraphics[width=0.99\linewidth]{Fig_01.jpg}
	\caption{Binders traditionally used for DMM. Source: modified after \cite{Ahnberg}}\label{fig1}
\end{figure} 

The tested specimens included soil of the following types, according to the Unified Soil Classification System (USCS): 1) Inorganic clays of low to medium plasticity, gravelly clays, sandy clays, silty clays, lean clays (CL); 2) Inorganic silts and very fine sands, rock flour, silty \hl{or} clayey fine sands or clayey silts with slight plasticity (ML). The specimens belong to the category of fine-grained soil \hl{consisting} of silts and clays with liquid limit less than 50\%. 

For low plasticity silty clay (CL), OPC has the best effects, followed by the lime-OPC blend and lime with equally good effects. Likewise, high plasticity clay (CH) is best \hl{stabilised} by the OPC with excellent effects, followed by the lime-OPC mixture and lime with \hl{an }equally good effect. The experimental results from the existing works show \hl{a} certain difference between the stabilisation and fieldwork using DMM, since the last one is used in the \emph{in-situ} conditions with dominating soft \hl{high-plasticity} clay (CH). Therefore, such methods does not include laboratory-based compaction works. At the same time, soil compaction is one of the essential parameters required for \hl{stabilisation}. 

Following different national standards and approaches on soil \hl{stabilisation} using one binder, the recommendations exist \hl{for} laboratory workflow \citep{ArifulHasan,Rogersetal2006}. According to these references, traditional and modified techniques of \hl{stabilisation} provide examples regarding the improving or optimising the existing methods \cite{Roduneretal2015,Kichouetal2015,Wangetal2020,Fahoum,Lindh202210,Lemenkova202222,Sasser,Lindh2003,Sisakhtetal2015}. For the assessment of binder blends, multiple and advanced experimental designs are preferable, including those adopted from different disciplines or domains. Thus, the test for the Initial Consumption of Lime (ICL) evaluates the amounts of lime which should be added in order to ensure that the required pH level is reached. The lower level guarantees that the effective pozzolanic reactions of soil with lime will occur. 

\begin{figure}[H]\centering
	\includegraphics[width=0.95\linewidth]{Fig_02.jpg}
	\caption{Boundaries best suited for various binders. Source: modified after \cite{Assarson}}\label{fig2}
\end{figure}   
\unskip

In contrast,\hl{ }if \hl{the} binder consists of the lime-GGBFS blend, there is no need to achieve the pH to ensure strength development, because lime and GGBFS can lead to this effect. With regard to the above\hl{,} complex design experiments on soil \hl{stabilisation} can be classified in \hl{two} groups: 1) to assess the amount of required \hl{stabilising} agents; 2) to test the ratios of binders with respect to the proportions of the selected components in a blended mixture. 

%\newpage
%\newpage\phantom{} 
\section{Statistical approach}
Statistical analysis ensures quality improvement, optimisation\hl{, }and economisation of the \hl{soil stabilisation workflow}, which is otherwise a very expensive, time-consuming and laborious process. Despite the similarities in\hl{ high-level }standard tests on binder selection, the advanced approaches using statistical methods of data analysis are fundamentally different in the following aspects: 1) statistical analysis enables using\hl{ }contextual information for a more robust evaluation of binder proportions; 2) statistical analysis is capable \hl{of} leveraging arbitrary features of soil collected in real \emph{in-situ} conditions which may vary considerably; 3) statistical analysis is embedding the standard context of soil classification data in the local features of soil. 

Statistical analysis enables to separate local features and standard classification types of soil which are matched and kept separately for optimisation of binder blends using their different matrices. As a result, a number of statistical methods performed \hl{a} verification of soil \hl{stabilisation} performance using integration of geotechnical engineering approaches and computational analysis \citep{Eissaetal2022,ORourke,Najjar,Kallenetal2016,Ojo,Lindh2021a}. Following this experience, the technique of the Response Surface Methodology (RSM) is adopted from the industrial areas including the automotive, chemical\hl{, }and process industries. The RSM proved to be an effective method to improve the reliability of binder testing. Besides, it helps the exploitation of new processes, optimises performance\hl{, }and improves the experimental design, as demonstrated in selected relevant studies \citep{SUBHA2017285,BIAN2021116745,Myres,CHEN2022125723}.

\section{Stabilisation quality control}

Robust estimation methods used in geotechnical engineering such as soil stabilisation are performed under real conditions which may include outliers in statistical data analysis. For instance, \emph{in-situ} setting may vary locally and include various soil grading, water content or mineral content in specimens. Therefore, tested methods should be fitted to the prevailing natural conditions that often include a large number of variables. To achieve the high level of soil strength, which indicates the quality of the \hl{stabilisation} process, blended mixtures should also be independent from variations in binder amount, homogeneity of mixtures\hl{, }and compaction. This ensures the objectivity of the statistical experiment using separated variables. Despite the complexity of the adopted approaches, it is important to optimise the workflow of data analysis, in order to achieve the ultimate goal of optimisation of binders. This ensures the analysis of the performance of soil in various scenarios, e.g.\hl{ }with varied proportions and content of binders, water ratio or with testing various soil samples.

Defining the determinant parameters for the stabilised soil \hl{is} essential to design high-quality soil-binder blend. The existing methods of parametrisation \cite{LindhNL} are adopted in the experimental work using background knowledge and pre-processing tests. These steps are crucial for the smooth workflow of the laboratory experiments. With this regard, it is principally important to note that \hl{the} evaluation of the amount of binders differs from the procedure of testing their types. These are the two separate subproblems that are to be solved in the same framework\hl{,} but using different evaluation approaches, in order to match the optimisation criteria of soil \hl{stabilisation}. 

The core of the evaluation method is estimating the amount of binder\hl{, }which is required to define lower and upper limits of the binder in a mixture. It strongly depends on the soil type\hl{ }as well as its mineral and moisture content. The potential solutions to this include a series of\hl{ }pre-tests integrated with the empirical evaluation of soil samples. Such pre-processing can utilise a few specimens and 1-2 response variables. 

The final design of the presented work was performed based on the principles discussed above and considering the results of the data preprocessing. The difference in the effects \hl{of} binders on \hl{the} gain in soil strength \hl{allowed us} to evaluate their effectiveness on the quality of soil \hl{stabilisation}. The decision on \hl{the} selection of binders that fits best to the given soil was performed with regard to the following response variables: 1) OMC; 2) Compaction energy required for a specific water content, Moisture Condition Value (MCV); 3) Uniaxial Compressive Strength (UCS) at different curing periods; 4) Water content after mixing of soil with binders which was assessed to evaluate and compare the effects from different agents on binding water in a mixture.

The aim of the undertaken study was a quantitative-qualitative analysis of binders for soil \hl{stabilisation}. First, we examined the types of binders which suit best for \hl{stabilisation} of the given soil type, that is, clay till. Second, we optimised the amount of these binders required to achieve the desired quality of soil using methods of statistical analysis. As a results of these tests, strength characteristics of soil \hl{were} significantly improved upon completion of the stabilisation process. The maximum UCS was achieved with the optimum amount of binder. The statistical analysis \hl{allowed} to ensure the economic improvements of works through the optimisation of the process. The CCD and the BBD methods were applied and integrated for \hl{the} quantity evaluation of binders.

\section{Results}
\subsection{Central Composite Design (CCD) as a Factorial Experiment}
The method of the Central Composite Design (CCD) is one of the most accepted \nth{2}-order experimental designs in engineering practice. It is based on the response surface methodology (RSM) \cite{Montgomery2019} and matrix approach, as shown in Equation \ref{eq1}. 

\begin{equation}\label{eq1}
\begin{bmatrix}
\alpha & 0 & 0 & \ldots & \ldots & \ldots & 0 \\
-\alpha & 0 & 0 & \ldots & \ldots & \ldots & 0 \\
0 & \alpha & 0 & \ldots & \ldots & \ldots & 0 \\
0 & -\alpha & 0 & \ldots & \ldots & \ldots & 0 \\
\vdots & \ldots & \ldots & \ldots & \ldots & \ldots & \vdots \\
0 & 0 & 0 & \ldots & \ldots & \ldots & \alpha \\
0 & 0 & 0 & \ldots & \ldots & \ldots & -\alpha \\
\end{bmatrix}
\end{equation}

It comprises the $2^\emph{k}$ factorial with n\textsubscript{f} runs, 2*\emph{k} axial runs, and n\textsubscript{c} centre runs. The \emph{k}=2 CCD is presented in Figure \ref{fig3}. The mathematical approach is based on the matrix\hl{, }which is derived from the axial points, with 2*\emph{k} rows depending on the factors\hl{, }which signify \hl{the} binder content. In our case, each factor is a binder\hl{, }which is consecutively located at \emph{a} and all other factors which might influence soil \hl{stabilisation} are set \hl{to} zero. This enables to evaluate the influence of binders as factors controlling the gain of strength. The value of \emph{a} is opted in accordance to the runs of binder testing. Changing values in binders enables to achieve the required properties of soil in the design of \hl{the} mixture, which leads to the increased gain in soil strength.

The input of each factor in this approach is defined at least at three levels: low, high and medium. The response \hl{corresponding} to the general model is presented in Equation \ref{eq2}.

%\begin{fleqn}
\begin{equation}\label{eq2}
\begin{aligned}[b]
y=\beta_0+\beta_{1}x_{1}+\ldots+\beta_{\kappa}x_{\kappa}+\beta_{12}x_{1}x_{2}+\beta_{13}x_{1}x_3\\
+\ldots+\beta_{\kappa-1,\kappa}x_{\kappa-1}x_{\kappa}+\beta_{11}x_{1}^2+\ldots+ \\\beta_{\kappa\kappa}x_{kappa}^2+\varepsilon
\end{aligned}
\end{equation}
%\end{fleqn}

where $\beta_i$ represents the regression coefficient and $\varepsilon$ -- the statistical error, so that $\varepsilon\in N(0,\sigma^2)$.

As can be inferred from the formulation of the feature matching problem, fitting a model to the observed values of the dependent variable \emph{y} includes three essential steps of data processing: 

\begin{itemize}
	\item main impacts for factors as components of binder blends $x_1, \ldots x_\kappa$;
	\item interactions between the \hl{stabilising} agents in binder blends as factors $x_{1}x_{2}, x_{1}x_{3}, \ldots, x_{\kappa-1}x_{\kappa})$; 
	\item quadratic components of factors $x_{1}^2, \ldots, x_{kappa}^2$. 
\end{itemize}

Modelling binder blends also provides useful information for detection of shear strength development, which can be measured additionally. The presented models for mixture data included the two models used for testing binder blends: the quadratic model and the special cubic model. 

Respectively, the quadratic model is formulated using the Equation \ref{eq3}:

\begin{equation}\label{eq3}
\begin{aligned}[b]
y=\beta_{1}x_{1}+\beta_{2}x_{2}+\beta_{3}x_{3}+\beta_{1,2}x_{1}x_{2}+\beta_{1,3}x_{1}x_{3}+\\\beta_{2,3}x_{2}x_{3}+\varepsilon
\end{aligned}
\end{equation}

where $\beta_i$ represents the regression coefficient and $\varepsilon$ -- the statistical error, similarly to the Equation \ref{eq2}.

Correspondingly, special cubic model can be expressed as demonstrated in the Equation \ref{eq4}:

\begin{equation}\label{eq4}
\begin{aligned}[b]
y=\beta_{1}x_{1}+\beta_{2}x_{2}+\beta_{3}x_{3}+\beta_{1,2}x_{1}x_{2}+\beta_{1,3}x_{1}x_{3}+\\\beta_{2,3}x_{2}x_{3}+\beta_{1,2,3}x_{1}x_{2}x_{3}+\varepsilon
\end{aligned}
\end{equation}

where $\beta_i$ represents the regression coefficient and $\varepsilon$ represents the statistical error, respectively. 

As one can note, the difference between these models consists in a special approach of cubic model\hl{, }which applies a test for interactions between all the factors $x_{1}x_{2}x_{3}$, sf. Equations \ref{eq2}, \ref{eq3} and \ref{eq4}.

Evaluating the effects \hl{of} binders on soil \hl{stabilisation} through a series of experiments can be used to detect the influence of each component of blend on the particular soil specimen with regard to the soil type and specifics of its mineral and moisture content. With this regard, we have performed the quantitative evaluation of \hl{the} binder reaction and assessed the effectiveness of the low, medium and high levels of binder content independently. To ensure the robustness of the quantitative evaluation, a single level was chosen initially, i.e.\hl{ }2.5\% binder per dry weight of soil. 

The initial proportion of binder, i.e.\hl{ }2.5\%, by weight of dry soil was used to make a groups of soil samples \hl{stabilised} by single binders prior to blended mixtures where proportions were selected in varying ratios to obtain the optimum percentage of each additive. The tested experiments were applied for clayey expansive tills of types CL (clays of low to medium plasticity to lean clays) and ML (inorganic silts and tills with silty of clayey fine sands or clayey silts with slight plasticity). Other types of soil would require different binders, since the reaction of binders with soil particle differs, depending on the type of soil: coarse-grained, medium-grained or fine grained. Thus, clay and silts are best \hl{stabilised} by lime (quicklime or hydrated lime), while for sandy soils\hl{, }it is more effective to use cementitious binders or pure OPC.

When using blended binders \hl{consisting} of three different \hl{stabilising} agents, such as OPC, lime and GGBFS, the blends can be fabricated using the defined proportions of binders as follows:

\begin{itemize}
	\item [--] Blend 1: OPC = 30\%, lime = 50\%, GGBFS = 20\%
	\item [--] Blend 2: OPC = 50\%, lime = 50\%, GGBFS = 0\%
	\item [--] Blend 3: OPC = 100\%, lime = 0\%, GGBFS = 0\%
\end{itemize}

Using statistical and combinatorial methods CCD, BBD\hl{, }and SLD, we defined that adding binders\hl{, }in general\hl{, }positively affects the development of strength, with lime being the best binder for \hl{stabilisation} of expansive clayey tills, followed by slag (GGBFS) and cement (OPC). \hl{The tested} binders strengthened the strength \hl{of} the specimens of high clay tills that have a high plasticity with the best performance showed by lime, followed by slag (GGBFS)\hl{, }and cement (OPC). 

Such effects are explained by the high content of active components in these binders, such as aluminium (III) oxide, calcium oxide\hl{, }and active silicon dioxide (silica), which increase\hl{ }the binding properties between the particles of \hl{soil-binder} mixtures and strength. Compared to single binders, a mixture of lime, slag (GGBFS)\hl{, }and cement (OPC) according to the proportions in Blend 1 (OPC = 30\%, lime = 50\%, GGBFS = 20\%) demonstrated the best effects on \hl{stabilisation} of clayey tills. Note that the hypothesis is made concerning the impact of factors, i.e.\hl{ }binders, to analyse the effects from \hl{a} combination of values from these factors \cite{StatSoft}. Using this general model is beneficial in several aspects\hl{, }as follows:

\begin{enumerate}
	\item The CCD method is adaptable in search of optimum in the response from factors. It enables to find the extreme values and saddle points\hl{, }which correspond to the optimal regions.
	\item The CCD presents feasible solutions to evaluate the parameters.
	\item The application of the CCD model is possible for different domains of the industrial design with specific assumptions.
	\item Additionally, although the CCD model is not suitable for the whole factorial range, an iterative process of gradual approximation through the loops of the adjustments of the true response function can result in \hl{an} optimised operation region achieved in parts, e.g.\hl{ }a narrow range around the best values.
	\item A combinatorial approach employed for finding the decisions can be used auxiliary for determining the directions of optimum, based on the choice of levels of the independent factors of the model.
	\item The \nth{2}-order models and the CCD as a factorial experiment\hl{ }can be used complementarily. 
\end{enumerate}

\begin{figure}[H]
	\includegraphics[width=0.75\linewidth]{Fig_03.jpg}
	\caption{(a): CCD for k=2 with start points and rotatable \nth{2}-order designs. (b): $3^2$ factorial design; (c): A comparison between the 2-factor CCD (continuous line) and a $2^2$ factorial design (dashed line).}\label{fig3}
\end{figure} 

The CCD with $\kappa=2$ is a $2^2$ factorial design which employs the combinatorial approach to solving the optimisation problem with 4 axial runs, Figure \ref{fig3}(a). This design includes loops used to fit the \nth{2}-order model \cite{Montgomery2019}. \hl{In case of} a face-centred CCD with a $\kappa=2$, the treatment combination is applied exactly as in a $3^2$ design, Figure \ref{fig3}(b). The discovered correspondences by the CCD and factorial design also include certain constraints on factor combinations for a standard CCD. Thus, a high level of both binders (OPC and lime) is not applicable since both binders require adding water, which leads to the conflict of variables. Figure \ref{fig3}(c) demonstrates this difference in \hl{the} applicability region. The observed cases in the first experimental design comprise OPC, lime and water as independent variables in CCD for $\kappa=3$.   

\subsection{Specimen collection and processing}
The soil used in this study is presented by specimens of clay till collected in the Yttre Ringvägen, a ring road in Malmö in southern Scania, Sweden. This soil is being used as an embankment fill material for the connecting road for the Öresund bridge. The specimens were excavated from the test pit in the location Petersborg. During the laboratory workflow, all the materials were passed through a 19 mm sieve, to exclude coarse-grained specimens from further processing, i.e.\hl{ }soil with grain size over 19 mm. 

\begin{figure}[H]
	\includegraphics[width=0.99\linewidth]{Fig_04.jpg}
	\caption{Water content in soil used as raw specimens for stabilisation tests. (a): test 1; (b): test 2.}\label{fig4}
\end{figure}   

The processing included the evaluation of factors as response variables, which were the following\hl{:} moisture content value (MCV), density, unconfined compressive strength (UCS)\hl{, }and change in water content after mixing of soil with binders. The material was systematically mixed and placed in sealed containers\hl{, }which were stored in a climate room with a maintained temperature of 8\degree C and a constant humidity of 85\%. The maintenance of temperature at 8\degree C aimed at simulating curing conditions of the specimens close to the natural conditions of Swedish environment during cool \hl{periods}, which are a temperature of 8\degree C and \hl{a} humidity of 85\%. At these conditions, soil samples were stored in a curing chamber for 28 days of curing period. Besides, such temperature and humidity conditions maintain real case environment during the freezing cycles and air drying\hl{, }which enables stimulating the actions of binders. As a result, the reaction of soil-binder mixtures performs effectively with particles of binder filling the pores of soil, which ultimately improves the compressive strength of soil. The curing period of one week was followed by the preprocessing\hl{, }which was performed to set\hl{ }up the test design. Water content measured in \hl{the} stored soil used as raw material in stabilisation tests is demonstrated in Figure \ref{fig4}. 

Water content demonstrated variations in the test range with theoretical value at 14.8\%, modelled values ranging from 15.5 to 16.5\%, and actual water content in the clay till lying in a range \hl{from} 8.4\% to 11.2\% (set 1) and 13.2\% to 15.8\% (set 2), Figure \ref{fig4}. The variation in water content and grading was minimised during the experiment.

\subsection{Effects from OPC and quicklime}
The OPC and quicklime were tested as independent variables during the \nth{2} experimental design of the face centre CCD for $\kappa=2$. The face centre design was applied to compare the performance and the effect of both mixed and unmixed binders, \emph{i.e.}, OPC and quicklime. In this case, the design shown in Figure \ref{fig3}(a) was changed to that shown in Figure \ref{fig3}(b) with low values assigned to zero. Statistically, this conforms that 8 of the tested design points were comprised of the stabilised material and one \hl{unstabilised}, i.e., the \hl{unstabilised} specimens form the part of the empirically tested region. 

Here\hl{, }two design points correspond to the pure lime and pure OPC, respectively, for the tested specimens with low values assigned to zero. The meaning of this step is to avoid negative values in the star (target) points which would arise if the low values in a standard CCD were set to zero which would result in the infeasible design of the experiment. Figure \ref{fig5}(a) and Figure \ref{fig5}(b) show the results from these empirical tests with the difference that Figure \ref{fig5}(a) only depicts the significant effects shown in a surface, while Figure \ref{fig5}(b) illustrates all the effects which were considered to model the surface. 

The response surface in Figure \ref{fig5}(a) is represented by the fitted function using the following Equation \ref{eq5}, (cf. Eq. \ref{eq1}). 

\begin{equation}\label{eq5}
\begin{aligned}[b]
w=15.2 -27*OPC-63.7*Ql+738.0*OPC*Ql 
\end{aligned}
\end{equation}

This equation only includes the significant terms as of OPC, lime and the interaction term OPC*quicklime (here: Ql). Correspondingly, the response surface in Figure \ref{fig5}(b) is expressed by the following fitted function in Equation \ref{eq6}:

\begin{equation}\label{eq6}
\begin{aligned}[b]
w=15.2-31.1OPC+68.4OPC2-31*Ql-\\1634.2*Ql+738*OPC*Ql
\end{aligned}
\end{equation}

In Equation \ref{eq6}, all the terms are included in the computing of the fitted function. The experimental loop comprised a series of tests with 40 specimens. This signifies that each point in the design signifies four specimens, except for the centre point\hl{, }which stands for eight samples. Major linear effects and the 2-way interconnections significantly impacted the variations in the water content. The interaction part of the Equation \ref{eq6} is positive, which means that using a mixture of the Ql and OPC did not reduce overall water content in the same way as\hl{ }using the pure stabilisation agent. This information can be used for the cases when soil specimens are sensitive to changes in water content, as in \hl{the} case with clay till from southwest Scania. 

\begin{figure}[H]
	\includegraphics[width=0.99\linewidth]{Fig_05.jpg}
	\caption{(a): The significant differences in water content depend on type and amount of binder; (b): The difference in water content depends on the type and amount of binder not adjusted for significance}\label{fig5}
\end{figure}   
\unskip

Moreover, \hl{stabilised} and \hl{unstabilised} soil behave distinctly regarding some parameters, such as OMC or grading. Therefore, it is a challenging task to include them in the same design. To support this through a comparative analysis, the relations between the MCV and water content in \hl{stabilised} and \hl{unstabilised} soil are demonstrated\hl{, }respectively\hl{, }in Figure \ref{fig7}. 

\subsection{ANOVA}
We estimated the design of the experiment by the analysis of variance (ANOVA) technique with the significance level at 0.05 for tested binders. For \hl{the} statistical evaluation, we used both pure and blended stabilisation agents\hl{, }which included quicklime\hl{, }hydrated lime, ordinary Portland cement (OPC)\hl{, }and their blends in various forms. The ANOVA was used to confirm that there \hl{are} no linear or quadratic main effects from these factors and no interaction between them, respectively. The results of the ANOVA are summarised in Table \ref{tab:2} with \emph{p} levels described as the probability of rejecting the null hypothesis when it is actually true. In the case of $p=0.05$, this means that there \hl{is a} 5\% probability that any relation between variables identified in the samples is random \cite{StatSoft}. The notations used in Table \ref{tab:2} are defined in the List of Symbols \ref{symb}.

Table \ref{tab:2} demonstrates that there is an interaction between the quicklime and the OPC, as indicated in Figures \ref{fig8} and \ref{fig9}. The interaction between the quicklime and the OPC is significant even at $p=0.01$. These results are also presented in the Pareto chart in Figure \ref{fig6} where a column presents the effects from the ANOVA summarised in Table \ref{tab:2}. These effects are organised and presented categorically, \hl{ordered} by the magnitude of values. The Pareto chart enables to estimate the magnitude, or the values of the effect, which should be achieved to be statistically significant according to the level \emph{p}. Thus, Figure \ref{fig6} shows that the OPC, Ql and the interaction between\hl{ }are significant at \hl{a} level of $p=0.05$. The quadratic effects from the OPC and Ql are not significant. 

\begin{figure}[H]
	\includegraphics[width=0.99\linewidth]{Fig_06.jpg}
	\caption{Pareto chart of the effects that notably change the variable W\_AFTER at p=0.05.}\label{fig6}
\end{figure}   
\unskip

%\begin{table}[htbp]
\begin{table*}[t]
\centering
\caption{ANOVA statistical analysis. Dependent variable: water content upon stabilization (W\_After) \label{tab:2}}
\begin{tabularx}{\textwidth}{c *{6}{Y}}
\toprule
Factor & \multicolumn{5}{l}{Anova; $R^2=0.79783$; Adj:0.76809, 2-factors, 1 block, 40 Runs, MS‑res=0.10063. } \\
\cmidrule(lr){2-4} \cmidrule(l){5-6}
  & SS & df & MS & F & p \\
\midrule
(1) Cem. (Linear)  & 8.3162 & 1 & 8.3162 & 82.640 & 0.0000 \\
Cem. (Quadratic) & 0.0353 & 1 & 0.0353 & 0.351 & 0.5572 \\
(2) Q.lime. (Linear) & 4.1409 & 1 & 4.1409 & 41.148 & 0.0000 \\
Q.lime (Quadratic) & 0.2492 & 1 & 0.2492 & 2.476 & 0.1247 \\
1(L) by 2(L) (interaction) & 0.7843 & 1 & 0.7843 & 7.793 & 0.0085 \\
Error & 3.4215 & 34 & 0.1006 & -- & -- \\
Total SS & 16.923 & 39 & -- & -- & -- \\
\bottomrule
\end{tabularx}
\end{table*}

There are no \nth{2} order terms in this equation, since it is reduced and modified according to the values of the significant effects, cf. with Table \ref{tab:2}. All the data points represent the same type of soil for \hl{a} robust statistical experiment. Although there is a difference between the reaction from various stabilisation agents with soil, all the binders fit the same regression line\hl{,} which has an R2 of 0.7759 for stabilised soil, based on 72 observations tested in this study. For the \hl{unstabilised} soil, the R2 of the regression line is 0.9303 and the fit is based on 18 observations\hl{.} 

\begin{figure}[H]
	\includegraphics[width=0.99\linewidth]{Fig_07.jpg}
	\caption{MCV vs water for treated and raw soil. The 95\%-confidence bound is shown by dash-dot lines.}\label{fig7}
\end{figure}   
\unskip

Figure \ref{fig7} illustrates that the MCV for the stabilized soil specimens becomes less sensitive to the variations in water content compared to the unstabilized raw soil before curing. From this comparison one can conclude that using different types and different amounts of binders in the same CCD concerning the response variable MCV is possible and results in effective soil stabilization. If a pure binder should be a part of a Box-Behnken Design, it is considered instead of using a face center CCD approach. We base the recommendation on the results from four other statistical runs, in which the standard CCD and the face-centred CCD were compared and analysed, respectively. 

The Pareto charts of the standardised effects from four different runs are shown in Figure \ref{fig8}. Here two different specimens of clay till were evaluated by tests and modelled using the response variable approach with regard to the changed water content before and after stabilisation ($\Delta W$). Respectively, the results are shown in the subfigures 1 and 3 of Figure \ref{fig8}, which represents these samples with a face-centred CCD. 

The subplot 1 in Figure \ref{fig8} shows that there are three significant linear effects from cement (L), the quadratic effect from cement (Q) and the linear effect from lime (L). Both linear effects from the OPC and those of quicklime increase $\Delta W$, while the quadratic effect from the OPC decreases $\Delta W$. This is owing to the disturbance from the unstabilized specimens which are the part of the design experiment. The results from the subplot 1 in Figure \ref{fig8} can be compared with those in the subplot 2 showing two significant linear effects from the OPC and Ql, respectively. 

The results shown that the quadratic effect of OPC is not significant at $p=0.05$ in a standard CCD test. The difference in the results is explained by the fact that in a face-centred CCD the unstabilized soil affect the final results. The subplot 3 in Figure \ref{fig8} shows the significant linear effects of Ql. The subplot 4 in Figure \ref{fig8} shows the significant linear effects from Ql and OPC. Adding 2.5\% of Ql to the soil reduced water content, which means that the results in the subplot 3 are affected by varied water content. In order to check the presents of the outliers in a general data pool, the predicted values were plotted against the residual ones, Figure \ref{fig9}, which demonstrated usual data pattern confirming the robustness of the experiment. 

\begin{figure*}[h!]
	\includegraphics[width=0.99\linewidth]{Fig_08.jpg}
	\caption{Pareto charts: modelled responses from four different tests. Subfigures 1 and 3 show face centre CCD, while 2 and 4 mean standard CCD. The response variable is modified water content before and after stabilization}\label{fig8}
\end{figure*}   
\unskip

\begin{figure*}[h!]
	\includegraphics[width=0.99\linewidth]{Fig_09.jpg}
	\caption{Predicted against residual values from the tests}\label{fig9}
\end{figure*}   
\unskip

\subsection{Box-Behnken Design (BBD)}
The Box-Behnken Design is an experimental design which determines the amount of binders \cite{Boxetal2005}, originally developed for RSM by G. E. P. Box and D. Behnken in 1960. It fits a quadratic model with one containing squared terms, products of two factors, linear terms and an intercept. In the BBD, each factor is evaluated at the three hierarchical levels: low medium and high, which are placed at one of three equally spaced values.

By visual examination of the BBD in Figure \ref{fig10}, one can note that the origin is located at a cube corner and the points in the corners of the cube and face points are absent. Such a design of BBD quantifies the predicting of the response for a pure binder which does not includes the stabilisation agents, as it is constrained with the low and high levels. The benefits of the BBD design consist in its optimisation with respect to the economical limitations, because it is particularly useful in modelling very expensive tests and experimental runs of large datasets. 

In the BBD design the minimum number of factors is three. Thus, if the design has >3 various binders, the high level of all binders is not considered, which presents a positive modelling effect. This is because if the region of interest is between 1-5\% for the 3 binders, the total amount of binders will otherwise be 15\% if all high levels are used at the same time. However, this is not appropriate in most of the cases. Thus, the BBD designed is optimised for the minimum number of factors at 3. Therefore, the BBD design was applied to estimate the time of performance and effects from 3 different binders at 3 different curing periods before compaction. We used the BBD in case where the standard CCD could not be used due to the target (star) points with time variable was set to 1, 3 and 5 hours of curing time, respectively.

\subsection{Simplex Lattice Design}
Similar to the case of BBD, each factor is evaluated at the three hierarchical levels: low, medium and high. We used the proposed simplex design which consists of the two parts: 1) simplex-lattice; 2) simplex-centroid design. These two types of the design are complemented with constraints. The real-time conditions accurately record changing state of the computed variables. Thus, if binder requires an activator, e.g., GGBFS, its pure blend will affect the results of the test which requires the assignment of the upper constraint. Figure \ref{fig11} demonstrates the design for the three components and simplex lattice design of the polynomial order 2. 

\begin{figure}[H]
	\includegraphics[width=0.90\linewidth]{Fig_10.jpg}
	\caption{BBD for 3 factors with a centre point}\label{fig10}
\end{figure}   
\unskip

The existing constrains in a simplex design is as follows. The cycles of the tests occur on the boundary of the model and the predictions in the interior can be uncertain. Besides, the design cannot be used as a universal approach for testing binders, since it does not contain the complete mixtures, but only binary or pure mixtures. At the same time, the interactions between all the binder components as stabilizing agents might be negative as well. The solution in such cases consists in adding the interior points, in case if the complete design region is of interest and requires extrapolation. 

To improve resolution, the experimental setup was adjusted by the repeated cycles with iterations increased from 6 to 10 following the existing methods \cite{electronics11223726}. The results regarding the fabrication of soil-binders mixture by simple lattice approach are presented in Figure \ref{fig12}. The subfigures A and B show the results from the augmented simplex lattice design with 20 tests, while subfigures C and D illustrate the result from a $3^2$ simplex lattice design with 12 runs.

\begin{figure}[H]
	\includegraphics[width=0.90\linewidth]{Fig_11.jpg}
	\caption{Simplex-Lattice design $3^2$}\label{fig11}
\end{figure}   
\unskip

\begin{figure*}[h!]
	\includegraphics[width=0.90\linewidth]{Fig_12.jpg}
	\caption{Pareto charts and RMS from simplex lattice design with UCS as a response variable. Subfigures C and D -- quadratic model; A and B -- \emph{idem} with interior points. Notable effects are visible in subfigures A and C}\label{fig12}
\end{figure*}   
\unskip

The subplots A and C show the estimated significant effects, while those presented in the subplots B and D are shown as a response surface. The interior points in this case did not affect the final results. Empirically, both the response surfaces demonstrate a positive interaction between lime and GGBFS regarding the gain in UCS of soil. Based on the results of this study, we recommend to avoid normal factor designs in case if single stabilizing agents are selected, because zero value of normal factor includes raw specimens. In case it exceeds zero value, only blended binders should be tested. Therefore, our method can be adopted if the goal of the test is to adjust the most suitable content of binders for soil stabilization, including specimens collected from other regions. We evaluated the mixture design of the ternary statistical modelling applied on mixtures of binders for soil stabilization, which proved to be effective and robust approach. Nevertheless, the type of soils and variation in water content in sample specimens should always be taken in the consideration for geotechnical engineering, prior to construction works. 

The performed experimental design proved to be a perfect tool for improving the efficiency and quality of various methods of testing soil quality and performance after stabilization with various binders: OPC, slag, and lime tested for Swedish soil. The soil material included clay tills collected from the embankment fill from the test places connecting road for the Öresund bridge, located in southern Scania region of Sweden. Besides the comparisons of binders and the performance of soil after curing we conducted the statistically significant test ANOVA implemented in Statistica. We shown that the difference between the testing methods is statistically significant as \emph{p}-values are all smaller than 0.05 significance level. Namely, we noted that OPC, lime and the interaction between them are significant at $p = 0.05$ and that the quadratic effects of the OPC and lime are not significant.

\section{Discussion}
In this study, we proposed an approach to evaluation of the amount and types of binders necessary for soil stabilization. In some cases, the amount of binder is set to a defined permanent value, or the interaction between various stabilizing agents for a fixed level which is regarded as a mixture of designs that needs to be applied. Careful assessment of binder proportions enables to take the advantages from the stabilisation agents constituting binder blend. 

To enhance strength development during soil stabilization, we modelled the modifications of binders and tested their effects on a soil-binder mixture. The sum of the components in a mixture design is set as constant, i.e., the amount of binders is not an independent variable but always counts for 100\%, regardless of the ratio of binders in each particular test case. We presented a new framework for improvement of engineering properties of soil using advanced method of stabilization with blended binders. The effects of binder ratios and water content on soil strength were evaluated by statistical methods and Pareto charts.

Among the proposed mixed binder blends for various types of soil, tested Blend 1 is an example of the complete mixture, fabricated of all the three binders. Blend 2 is a modified binary blend where we excluded the influence of the GGBFS, and the impact of OPC and lime is equal (50/50\%). The single-binder Blend 3 is a pure binder containing OPC only up to 100\%. 

\section{Conclusions}
In this paper, we introduced a framework for selecting optimal amounts of binders used for soil stabilization with a case study of clay soils collected from the region of Yttre Ringvägen, which is a ring road in Malmö located in southern Sweden. A novel framework for adjusting binder content was proposed. We tested three approaches of the RSM techniques for evaluation of soil performance stabilised using various binder types. The results were compared and presented using Pareto charts. We used the following combinatorial statistical methods for data analysis and assessment: 1) Central Composite Design (CCD); 2) Box-Behnken Design (BBD); 3) Simplex Lattice Design. 

Given the large amount of data and processed materials in the constructions industry where tons of soil should be processed for evaluation to ensure the safety of the foundations prepared for constructed buildings, statistical analysis is one of the major challenges supporting and facilitating the laboratory experiments through data modelling. To deal with soil specimens, we propose incorporating statistical methods of CCD, BBD and SLD to refine the selection of binders prepared as components for blended mixtures. The increase of lime improved the stabilization performance followed by cement and slag. 

We devised a combinatorial method to use information on types and amount of 3 blends of binders for refining and optimising binder blends and then develop a modelling approach using Statistica software to use the refined formulae of binder blends as a practical application for soil stabilization. Blend 1, a mixture fabricated of the three binders shown the best results, followed by the Blend 2, which is a modified binary blend where the influence of slag is excluded, while the impact of OPC and lime is equal, and a single-binder Blend 3, a pure binder containing OPC. The experimental results show good improvement in binder blend matching using the approaches of CCD, BBD and SLD, which proves the significance of these methods for soil stabilization.

To demonstrate the capability of various binder and their influence on soil stabilization and strength development, we applied and tested two approaches: a design for evaluating the amount of binder and a design to evaluate the interaction between several binders. The proposed method is shown to be capable of incorporating the RMS, which as an effective tool for geotechnical experiments. We demonstrated that the CCD using robust functionality is essential for the assessment of the amount and type of the stabilizing agents necessary for effective soil stabilization. Additionally, we indicated that BBD can be used for statistical tests in civil engineering and geotechnical tasks. The statistical analysis included quadratic model, regression analysis and ANOVA test to analyse the efficiency of measurements.

We recommend to take into account the following closing remarks for similar related works:
\begin{itemize}
	\item The CCD is effective for blended binder with 2 agents for modelling the content of mixture and ratio in between the components.
	\item The CCD and BBD techniques both can be applied for blends with >3 stabilizing agents to optimise their amounts in each case.
	\item A mixture design is effective to test the effects from the existing binders or adjust the new ones in cases of complex mixtures of binders (>3 elements of stabilizing agents in a mixture).
	\item The external factors (water content or temperature) can be processed as covariates in a function.
	\item The repetitive measurements with >2 sample tests should be used in every type of the experiment.
	\item For correct comparison between the unstabilized and stabilised specimens, a design should include only stabilised samples followed by the reference samples for the unstabilized ones.
\end{itemize}

The specific value of our research is that the methods and techniques have been borrowed from other industrial areas such as the automotive sector, chemical engineering and process industry. Thus, we used the RSM as an advanced and efficient technique which can be used to optimise binder types and amount for soil stabilisation. Next, the Pareto charts were presented for modelling responses from different tests with changed types and amount of binders and estimate the effects of binder ratios and water content on soil strength. Using this approach, we conducted the extensive experiments on soil stabilisation using new binders and novel methods of soil stabilisation which required different testing and evaluation methods to benchmark the performance of the framework. 

The enhanced properties of binders with regard to modified and updated performance can be applied in civil engineering. Otherwise, using traditional ways of testing in soil stabilization, is an expensive, time-consuming and laborious task in construction works. Further experiments on testing binder blends by means of statistical analysis demonstrated their potential applications in civil engineering, applied material science and geotechnical studies. 

This paper is an important point in further development of these directions, presenting the combined methods of statistical analysis, theoretical logic, data science and approaches of combinatorics for practical tasks of soil stabilization in civil engineering. Finally, we demonstrated the applications of combinatorics for optimal combination of binders in soil stabilization.

\section{Abbreviations and notations}\label{symb}
\begin{nomenclature}
\item{Adj}{Adjusted $R^2$}
\item{ANOVA}{ANalysis Of VAriance}
\item{BBD}{Box-Behnken Design}
\item{CCD}{Central Composite Design}
\item{df}{Degrees of freedom}
\item{DMM}{Deep Mixing Method}
\item{F}{F test. $F=MS_{treatments}/MS_{e}$}
\item{GGBFS}{Ground Granulated Blast Furnace Slag}
\item{ICL}{Initial Consumption of Lime}
\item{MCV}{Moisture Compaction Value}
\item{MS}{Means of squares}
\item{OMC}{Optimum Moisture Content}
\item{OPC}{Ordinary Portland Cement}
\item{R-sqr}{Coefficient of multiple determination (=$1-SS_{e}/SS_{t}$)}
\item{RSM}{Response Surface Methodology}
\item{SGI}{Swedish Geotechnical Institute}
\item{SS}{Sum of Squares, for each treatment, error and total}
\item{p}{p-level}
\item{UCS}{Unconfined/Uniaxial Compressive Strength}
\end{nomenclature}

\begin{acknowledgements}
We acknowledge the Cementa (HeidelbergCement AG Group -- Northern Europe), Nordkalk AB, Swedish National Board for Industrial and Technical Development (NUTEK) and Peab AB for their support of this study. The authors gratefully acknowledge  valuable assistance of Prof. Dr. Björn Holmqvist from the Department of Statistics, Lund University and Prof. Dr. Em. Igor Rychlik from the Centre for Mathematical Sciences, Chalmers University of Technology. We appreciate critical comments, corrections and remarks from the two anonymous reviewers which improved the initial version of this manuscript.
\end{acknowledgements}


\bibliographystyle{actapoly}
\bibliography{biblio}
%\nocite*{}

\end{document}
